Visualisations are essential for extracting insight from any dataset. Given the
scale of NaLoDuCo datasets, downloading them locally is impractical. Therefore,
visualisation methods must operate where the data resides—either in the cloud
or on institutional HPC systems.

We will develop visualisation functionality for both continuous datasets and
epoched datasets, where epochs are anchored around events identified by
advanced machine learning methods.

\subsubsubsection{Continuous visualisations}

Continuous visualizations will allow users to intuitively explore large-scale
behavioral and neural datasets spanning weeks to months. Users will be able to
zoom out for a high-level overview (e.g., a full month) or zoom in to examine
millisecond-level details, offering an interactive experience similar to Google
Maps for time series data.

To enable this, we will implement:

\begin{itemize}

	\item\textbf{Multi-Resolution Tiling:} Datasets will be preprocessed into
	tiles at various spatial and temporal resolutions, rendering only relevant
	tiles as users zoom.

	\item\textbf{Hierarchical Storage:} Data will be organized using formats
	such as Zarr or HDF5 to support chunked, multi-resolution access.

	\item\textbf{On-Demand Streaming:} Visualizations will stream data
	dynamically, adapting to the user's current view.

\end{itemize}

\subsubsubsection{Epoched and interactive visual analytics}

A key strength of our platform is its support for epoched visualisation and
interactive, closed-loop visual analytics, which together enable the discovery
and refinement of neural and behavioural patterns in long-duration datasets.

\paragraph{Machine Learning-Defined Events.} A core feature of our system will
be the ability to align epochs not just to experimenter-defined events, but
also to latent state transitions inferred via unsupervised methods (e.g.,
hidden Markov models, behavioural clustering, inverse reinforcement learning)

\paragraph{Closed-Loop Analytics.} There will be a closed-loop interaction
between visualisations and machine learning algorithms. In this loop:

\begin{itemize}

	\item \textbf{Machine learning algorithms} extract latent states, classify
	behaviours, infer structure, or forecast dynamics from NaLoDuCo data.

	\item These outputs feed into the visualisation engine to generate novel
	views (e.g., state-aligned rasters, dynamic embeddings, attention maps).

	\item \textbf{Users explore these visualisations interactively},
	discovering unexpected, task-agnostic, or contextual patterns.

	\item New queries and insights drive further rounds of machine learning
	analysis—closing the loop.

\end{itemize}

This design enables researchers to co-develop computational models and
scientific hypotheses iteratively, with human insight and machine inference
deeply intertwined.

\subsubsubsection{Deliverables}

\begin{enumerate}

	\item visualisations for continuous behavioural and neural recording

	\item visualisations for epoched behavioural and neural recording

	\item visualisations for model outputs

	\item indexing system to support intelligent visualisations

	\item deployment of the above items to allow users to visualise NaLoDuCo
	DANDI datasets on the cloud

\end{enumerate}
